% Generated by `taskgen exam render` from materials/sem2/exam/questions.yaml.
% Do not edit: change questions.yaml and run `taskgen exam render` (or `make exams`).
\examtitle[вопросов на экзамен]{Вопросы на экзамен}{июнь 2027}
\examliterature{hps,galbraith,crandall-pomerance,msu-ent,glukhov}

\topic{Дискретное логарифмирование и разложение на простые множители}
\begin{questiongroup}{1}{2}{\cite[Глава 2]{hps}}  % sem2:1
  \qpart{Задача дискретного логарифмирования в $\Zn{n}$. Применение DLP на практике: протокол Диффи"--~Хеллмана. Тривиальная оценка сложности DLP.}
  \qpart{Алгоритм Шенкса (Babystep-Giantstep), оценка сложности.}
\end{questiongroup}
\begin{questiongroup}{3}{4}{\cite[Глава 2]{hps}}  % sem2:3
  \qpart{Алгоритм решения DLP методом сведения к собственным подгруппам (по основанию, порядок которого "--- степень простого числа).}
  \qpart{Китайская теорема об остатках. Алгоритм Полига"--~Хеллмана, оценка сложности.}
\end{questiongroup}
\begin{questiongroup}{5}{7}{\cite[Глава 4]{hps} \cite{galbraith}}  % sem2:5
  \qpart{Теорема о коллизиях, её применение для оценки ожидаемой длины списков, с заданной вероятностью обладающих одинаковым элементом. Тривиальный алгоритм решения DLP с помощью поиска коллизии двух списков случайных степеней основания логарифма $g$.}
  \qpart{Понятие прямой орбиты элемента конечного множества, структура орбиты. Теорема о периоде циклической части орбиты.}
  \qpart{Решение DLP в $\Zn{n}$ с помощью $\rho$-алгоритма Полларда и алгоритма поиска цикла Флойда, оценка сложности.}
\end{questiongroup}
\question{8}{Решение DLP в $\Zn{n}$ с помощью $\rho$-алгоритма Полларда с выделенными точками, оценка сложности.}{\cite{galbraith}}  % sem2:8
\question{9}{Алгоритм решения DLP методом исчисления индексов (index calculus), оценка сложности.}{\cite[Раздел 3.8]{hps}, \cite[Раздел 6.4]{crandall-pomerance}}  % sem2:9
\question{10}{Основная теорема арифметики. Формулировка задачи разложения (факторизации) составного числа. Тривиальная оценка сложности задачи факторизации.}{любой учебник по теории чисел, например, Бухштаб, \cite{msu-ent}}  % sem2:10
\question{11}{Основной и расширенный алгоритм Евклида. Теорема Эйлера. Применение сложной задачи факторизации на практике: RSA.}{\cite{msu-ent}, \cite[Раздел 3.2]{hps}}  % sem2:11
\begin{questiongroup}{12}{14}{\cite[Разделы 3.6--3.7]{hps}, \cite[Глава 6]{crandall-pomerance}}  % sem2:12
  \qpart{Разложение на множители на основе разности квадратов. Трёхшаговая процедура разложения числа на множители.}
  \qpart{Понятие $B$-гладкого числа. Алгоритм Квадратичного решета.}
  \qpart{Оценка сложности и метод выбора параметров алгоритма Квадратичного решета (с выводом).}
\end{questiongroup}
\question{15}{Алгоритм факторизации $p - 1$ Полларда, оценка сложности.}{\cite[Раздел 3.5]{hps}, \cite[Раздел 5.4]{crandall-pomerance}}  % sem2:15
\question{16}{Критерий Эйлера существования квадратного корня по модулю простого числа. Решение квадратичного сравнения по модулю нечётного простого числа методом Чиполлы.}{\cite[Раздел 2.3--2.4]{glukhov} или Википедия}  % sem2:16
\question{17}{Лемма Гензеля, её приложение к решению квадратичного сравнения по модулю степени нечётного простого числа.}{\cite{msu-ent}}  % sem2:17

\topic{Эллиптические кривые}
\begin{questiongroup}{18}{20}{\cite[Разделы 6.1--6.2]{hps}, \cite[Глава 7]{crandall-pomerance}}  % sem2:18
  \qpart{Уравнение эллиптической кривой в форме Вейерштрасса. Определение эллиптической кривой, дискриминант.}
  \qpart{Операция сложения точек кривой, группа точек эллиптической кривой.}
  \qpart{Эллиптические кривые над конечными полями, конечная группа точек эллиптической кривой. Верхняя граница числа точек в группе, теорема Хассе (без доказательства).}
\end{questiongroup}
\question{21}{Быстрый алгоритм возведения точки в степень. Возведение в степень на основе тернарного разложения показателя. Оценка максимального и ожидаемого (среднего) количества элементарных групповых операций.}{\cite[Разделы 6.3--6.4]{hps}}  % sem2:21
\question{22}{Задача дискретного логарифмирования для группы точек на эллиптической кривой, ECDLP, тривиальная оценка её сложности. Применение ECDLP на практике: протокол Диффи"--~Хеллмана на эллиптических кривых.}{\cite[Разделы 6.3--6.4]{hps}}  % sem2:22
\question{23}{Решение ECDLP с помощью $\rho$-алгоритма Полларда, оценка сложности.}{\cite[Раздел 5.5]{hps}}  % sem2:23
\question{24}{Алгоритм Ленстры факторизации целого числа с помощью эллиптической кривой, его сложность (без доказательства).}{\cite[Раздел 6.6]{hps}}  % sem2:24

\topic{Решётки}
\question{25}{Решётка, базис решётки, определитель решётки. Задачи о кратчайшем и ближайшем векторе (SVP, CVP). Алгоритм Гаусса редукции базиса двумерной решётки, оценка сложности.}{\cite[Разделы 7.4--7.5, 7.13.1]{hps}, \cite[Раздел 17.1]{galbraith}}  % sem2:25 (new)
\question{26}{Ортогональность базиса, коэффициент Адамара. Решение задачи о ближайшем векторе алгоритмом Бабаи, его результат для «хорошего» и «плохого» базисов.}{\cite[Разделы 7.6, 7.13.3]{hps}, \cite[Разделы 18.1--18.2]{galbraith}}  % sem2:26 (new)
\question{27}{Ортогонализация Грама"--~Шмидта. LLL-приведённый базис (условие уменьшения размера и условие Ловаса). Алгоритм LLL, его корректность и оценка сложности.}{\cite[Раздел 7.13.2]{hps}, \cite[Глава 17]{galbraith}}  % sem2:27 (new)
\question{28}{Конгруэнтная криптосистема с открытым ключом. Атака на конгруэнтную криптосистему с помощью редукции базиса двумерной решётки.}{\cite[Разделы 7.1, 7.14.1]{hps}}  % sem2:28 (new)
